% !TEX root = ../master_thesis.tex
\chapter{Prompt Templates and Experimental Metadata}
\label{app:prompt-templates}

This appendix records prompts that materially affect the thesis methodology,
experimental results, benchmark protocol, or reproducibility. Routine prompts
used for coding, brainstorming, language polishing, or project management are
not included unless their outputs become part of the experimental protocol.

\section{Inclusion Rule}
\label{app:prompt-inclusion-rule}

A prompt is included when changing or omitting it could change a reported
metric, model comparison, benchmark construction step, annotation protocol, or
reproducibility claim. For each included prompt, the experiment records the
prompt text, model identifier, provider, date, key decoding parameters, input
modalities, output format, random seed if applicable, and source result file.

\section{Track A Pairwise Judgement Prompts}
\label{app:track-a-prompts}

\subsection{TA-PAIR-ZS-v1: Zero-shot Pairwise Visual Design Prompt}
\label{app:prompt-ta-pair-zs-v1}

\textbf{Purpose:} Baseline pairwise GUI screenshot judgement for Track A,
including the order-swapped reliability condition. The model receives two
screenshots labelled Image A and Image B and returns only the preferred letter.

\textbf{Metadata:}
\begin{itemize}
  \item \textbf{Track / task:} Track A, pairwise preference judgement
  \item \textbf{Prompt ID:} TA-PAIR-ZS-v1
  \item \textbf{Models used:} \texttt{gpt-4o} (OpenAI); \texttt{claude-sonnet-4-5-20250929} (chatanywhere-anthropic proxy)
  \item \textbf{Input modalities:} two PNG screenshots (Image A, Image B) + user-role text
  \item \textbf{Message structure:} single user message; no system message
  \item \textbf{Decoding parameters:} temperature = 0, max\_tokens = 5
  \item \textbf{Seed:} 42 (pair sampling); model temperature is 0 so output is deterministic
  \item \textbf{Order-swap:} applied; each pair is evaluated in both AB and BA orders
  \item \textbf{Dataset:} \texttt{biglab/uiclip\_human\_data\_hf} (HuggingFace), train split
  \item \textbf{Source:} \texttt{scripts/track\_a\_eval.py}, lines 44--52 (\texttt{ZEROSHOT\_PROMPT})
  \item \textbf{First use:} 2026-05-06 (gpt-4o pilot); 2026-05-19 (claude-sonnet-4-5 run)
  \item \textbf{Result files:} \texttt{scripts/results/track\_a\_eval\_gpt-4o\_20260506\_*}; \texttt{scripts/results/track\_a\_eval\_chatanywhere-anthropic\_claude-sonnet-4-5-20250929\_zeroshot\_20260519\_153820.json}
\end{itemize}

\begingroup
\scriptsize
\begin{verbatim}
You are evaluating two user interface screenshots for visual design quality.

Image A is shown first, Image B is shown second.

Which interface has better visual design overall? Consider layout, visual
hierarchy, clarity, and aesthetic quality.

Reply with ONLY the letter A or B. Nothing else.
\end{verbatim}
\endgroup

\subsection{TA-PAIR-RUBRIC-v1: Rubric-guided Pairwise Visual Design Prompt}
\label{app:prompt-ta-pair-rubric-v1}

\textbf{Purpose:} Rubric-guided Track A pairwise GUI screenshot judgement.
Adds four explicit design dimensions before requiring the same A/B output
format as the zero-shot condition. Used in the prompting-strategy ablation.

\textbf{Metadata:}
\begin{itemize}
  \item \textbf{Track / task:} Track A, pairwise preference judgement (rubric condition)
  \item \textbf{Prompt ID:} TA-PAIR-RUBRIC-v1
  \item \textbf{Models used:} \texttt{gpt-4o} (OpenAI)
  \item \textbf{Input modalities:} two PNG screenshots (Image A, Image B) + user-role text
  \item \textbf{Message structure:} single user message; no system message
  \item \textbf{Decoding parameters:} temperature = 0, max\_tokens = 5
  \item \textbf{Seed:} 42 (pair sampling)
  \item \textbf{Order-swap:} applied
  \item \textbf{Dataset:} \texttt{biglab/uiclip\_human\_data\_hf} (HuggingFace), train split
  \item \textbf{Source:} \texttt{scripts/track\_a\_eval.py}, lines 54--68 (\texttt{RUBRIC\_PROMPT})
  \item \textbf{First use:} 2026-05-19
  \item \textbf{Result files:} \texttt{scripts/results/track\_a\_eval\_gpt-4o\_20260506\_142748\_pair\_review/}
\end{itemize}

\begingroup
\scriptsize
\begin{verbatim}
You are an expert UI design reviewer evaluating two user interface screenshots.
Image A is shown first, Image B is shown second.

Score each interface on the following four dimensions before deciding:

1. Visual structure - layout, alignment, spacing, and visual hierarchy.
2. Information clarity - readability, labelling, and grouping of related items.
3. Aesthetic quality - typography, colour use, and overall polish.
4. Inferred usability - how easy the interface looks to use at a glance.

Briefly compare the two interfaces dimension by dimension in your head, then
pick the one that is better overall.

Reply with ONLY the letter A or B on the last line. Nothing else.
\end{verbatim}
\endgroup

\section{Track B GUI Generation Prompts}
\label{app:track-b-generation-prompts}

\subsection{Track B Generation Prompt Version History}
\label{app:track-b-generation-version-history}

Track B generation prompts changed repeatedly during smoke testing. To keep
the appendix focused on reproducibility rather than serving as a full prompt
archive, this appendix reproduces in full the frozen baseline prompt that
produces the reported Track B pilot and leaderboard artifacts,
\texttt{TB-GEN-v9}, the later F10 diagnostic prompts
\texttt{TB-GEN-v14} and \texttt{TB-GEN-v15}, and the compact-input candidate
\texttt{TB-GEN-v16}. All earlier smoke-test variants
(\texttt{TB-GEN-v1} through \texttt{v8}, \texttt{v12}, and \texttt{v13}) were
diagnostic iterations that did not contribute reported metrics; they are
summarized in Table~\ref{tab:track-b-generation-prompt-history} and remain
recoverable in full from the source template in
\texttt{scripts/generate\_track\_b\_ui.py}, per-run
\texttt{generation\_metadata.json} files, and the revision log when needed to
interpret a specific pilot artifact.

\begin{table}[ht]
\centering
\footnotesize
\begin{tabularx}{\linewidth}{>{\raggedright\arraybackslash}p{0.15\linewidth}>{\raggedright\arraybackslash}p{0.16\linewidth}>{\raggedright\arraybackslash}X>{\raggedright\arraybackslash}p{0.16\linewidth}}
\toprule
\textbf{Prompt ID} & \textbf{Use} & \textbf{Main change} & \textbf{Status in appendix} \\
\midrule
\texttt{TB-GEN-v1} & First smoke run & Baseline single-file HTML generation prompt for normalized Vision2Web pilot items. & Summarized only \\
\texttt{TB-GEN-v2} & Smoke iteration & Added explicit completion, same-file JavaScript handlers, workflow-action non-placeholder requirements, and responsive-layout constraints. & Summarized only \\
\texttt{TB-GEN-v3--v4} & Intermediate iteration & Refined workflow route handling before the fixed route-contract version. & Summarized only \\
\texttt{TB-GEN-v5} & Route-contract smoke run & Added a fixed route contract, a minimal route handler, no-scroll hash replacement, and stricter mobile overflow safeguards. & Summarized only \\
\texttt{TB-GEN-v6} & Reproducible smoke template & Made desktop workflow fidelity the hard generation target and clarified handling of unavailable downloadable assets. & Summarized only \\
\texttt{TB-GEN-v7} & Compact same-prompt comparison template & Keeps the same route/workflow contract as v6 but instructs models to produce a shorter complete implementation to reduce truncation and gateway-timeout risk. & Summarized only \\
\texttt{TB-GEN-v8} & Compact gated comparison template & Adds explicit exact workflow-label-to-route mappings and forbids inert elements for workflow-required controls after v7 produced an inert Local Forms control. & Summarized only \\
\texttt{TB-GEN-v9} & Per-item workflow-label prompt (frozen baseline) & Generalizes v8 by automatically extracting workflow click labels from each item's workflow file instead of hard-coding the F09 label-to-route list. Frozen as the practical baseline for the cross-item pilot batch. & Full text below \\
\texttt{TB-GEN-v12} & Machine-checkable workflow-control contract & Adds per-item \texttt{visible\_text}, \texttt{route\_id}, and \texttt{required\_element} entries to reduce ambiguity around exact clickable labels; did not resolve the exact-label case and was not adopted. & Summarized only \\
\texttt{TB-GEN-v13} & Required workflow-control navigation block & Replaces per-control examples with a complete navigation block to copy into the page header; both smoke attempts failed at the provider layer (HTTP 500) and produced no artifact. & Summarized only \\
\texttt{TB-GEN-v14} & F10 visual-grounding diagnostic & Removes the strong compactness bias, requires local images, and uses per-control \texttt{required\_element} entries while clarifying that workflow context phrases are not necessarily clickable controls. & Full text below \\
\texttt{TB-GEN-v15} & F10 control-contract diagnostic & Tightens v14 after F10 v14 failed: forbids non-contract \texttt{showPage()} routes and forbids replacing required control text with labels such as ``View Recipe''. & Full text below \\
\texttt{TB-GEN-v16} & Compact-input candidate & Keeps v15's exact-control contract but uses a compact route/evidence workflow summary, separates semantic visual targets from exact visible text controls, adds bounded generation, and uses deterministic ranked resources plus prototype-image downscaling in the generator input policy. & Full text below \\
\bottomrule
\end{tabularx}
\caption{Summary of Track B GUI generation prompt versions. Only the current
templates are reproduced in full in this appendix; prior versions are retained
through experiment artifacts and source history.}
\label{tab:track-b-generation-prompt-history}
\end{table}

\subsection{TB-GEN-v9: Per-Item Workflow-Label Generation Prompt}
\label{app:prompt-tb-gen-v9}

\textbf{Purpose:} Practical Track B baseline prompt for the cross-item smoke
batch. This version generalizes the v8 workflow-control idea by filling the
workflow label list automatically from each benchmark item's workflow file
rather than using a hand-written item-specific mapping.

\textbf{Metadata:}
\begin{itemize}
  \item \textbf{Track / task:} Track B, GUI generation
  \item \textbf{Prompt ID:} TB-GEN-v9
  \item \textbf{Version change:} revises TB-GEN-v8 by replacing the
  hard-coded F09 workflow-control list with an automatically extracted
  \texttt{\{workflow\_labels\}} slot
  \item \textbf{Models used:} \texttt{qwen3.6-35b-a3b} via GWDG/SAIA
  OpenAI-compatible API
  \item \textbf{Input modalities:} normalized requirement text, Vision2Web
  workflow JSON, local resource-path list, route contract, and prototype
  screenshots
  \item \textbf{Message structure:} system message plus one user message
  containing text and prototype images
  \item \textbf{Decoding parameters:} temperature = 0.2, max\_tokens = 20000
  in smoke tests
  \item \textbf{Dataset:} \texttt{zai-org/Vision2Web}, reduced Level 1 /
  Level 2 pilot subset
  \item \textbf{Source:} \texttt{scripts/generate\_track\_b\_ui.py}
  \item \textbf{First use:} 2026-05-26, Europe/Berlin
  \item \textbf{Result files:} \texttt{data/track\_b/items/F01\_1daycloud/generated/gwdg\_qwen36\_35b\_v9\_smoke/};
  \texttt{data/track\_b/items/F03\_about\_gitlab/generated/gwdg\_qwen36\_35b\_v9\_smoke/};
  \texttt{data/track\_b/items/F10\_gourmania/generated/gwdg\_qwen36\_35b\_v9\_smoke/}
\end{itemize}

\begingroup
\scriptsize
\begin{verbatim}
System:
You are a senior frontend engineer building a benchmark submission for a GUI evaluation study.
Create a polished, working website implementation that can be opened directly from a local index.html file.

User:
Prompt ID: TB-GEN-v9

Build a minimal complete single-file HTML implementation for the Track B GUI item below.

Hard requirements:
- Return exactly one complete HTML document ending with </html>.
- Include inline CSS and JavaScript only. No build step, backend, network access, external CDN, or markdown.
- Local assets must be referenced as ../../resources/<subpath>.
- Use the prototype screenshots for overall structure and visual cues, but produce a minimal benchmark implementation.
- Primary viewport is desktop: 1365px and 1920px. Mobile quality is not a hard gate.
- Implement every route in the route contract as a <section id="..." data-track-route>.
- Include window.__TRACK_B_ROUTES exactly as given and a showPage(routeId) handler before </body>.

Mandatory workflow controls:
- The following exact workflow click labels were extracted from the workflow checks. Every label below must appear as visible text inside an <a> or <button> element:
{workflow_labels}
- Never use <div>, <span>, <li>, or other inert elements for workflow-required click labels.
- Every workflow control must include data-route-target="<route_id>" and onclick="showPage('<route_id>')".
- If a click target is a feature phrase such as "GitLab Duo", use that feature phrase as the visible label.
- Choose the route target from the route contract that best matches the label and workflow validation.
- Secondary navigation items may be visible non-functional placeholders only if they are not listed in the extracted workflow click labels.

Strict compactness limits:
- Use one shared header and one shared footer only.
- For each route, use at most one h1, one short paragraph, one compact card/list, and one optional tiny table.
- Use at most 3 cards per route, 3 bullets per list, and 3 rows per table.
- Do not copy long requirement text. Use short labels and short summaries.
- Avoid decorative animations, large CSS frameworks, inline SVG art, exhaustive menus, and long repeated sections.
- If output budget feels tight, remove visual detail first. Never omit routes, workflow controls, JavaScript, or closing tags.

Route contract:
{route_contract}

Required JavaScript:
<script>
window.__TRACK_B_ROUTES = {route_ids_json};
function showPage(routeId) {
  if (!window.__TRACK_B_ROUTES.includes(routeId)) return;
  document.querySelectorAll('[data-track-route]').forEach(function(section) {
    section.hidden = section.id !== routeId;
  });
  document.querySelectorAll('[data-route-target]').forEach(function(link) {
    link.setAttribute('aria-current', link.dataset.routeTarget === routeId ? 'page' : 'false');
  });
  if (history.replaceState) history.replaceState(null, '', '#' + routeId);
  setTimeout(function() { window.scrollTo(0, 0); }, 0);
}
document.addEventListener('DOMContentLoaded', function() {
  var initial = location.hash ? location.hash.slice(1) : window.__TRACK_B_ROUTES[0];
  showPage(window.__TRACK_B_ROUTES.includes(initial) ? initial : window.__TRACK_B_ROUTES[0]);
});
</script>

Item ID: {item_id}
Source task: {source_task_name}
Source level: {source_level}
Use for dynamic validation: {use_for_dynamic}

Normalized requirement:
{requirement}

Workflow checks:
{workflow}

Available local resource paths:
{resource_paths}

Prototype screenshots are attached after this text in this order:
{prototype_list}
\end{verbatim}
\endgroup

\subsection{TB-GEN-v14: Visual-Grounded Workflow-Control Diagnostic Prompt}
\label{app:prompt-tb-gen-v14}

\textbf{Purpose:} F10 diagnostic prompt introduced after TB-GEN-v9 and cheap
second-model smoke tests produced bare or weakly interactive pages. This
version tests whether stronger visual-grounding instructions and an explicit
machine-checkable workflow-control contract improve generated artifacts without
changing the underlying F10 benchmark item.

\textbf{Metadata:}
\begin{itemize}
  \item \textbf{Track / task:} Track B, GUI generation
  \item \textbf{Prompt ID:} TB-GEN-v14
  \item \textbf{Version change:} revises TB-GEN-v9/v12 by removing strong
  compactness limits, requiring local images, and clarifying that workflow
  context phrases are not necessarily clickable controls
  \item \textbf{Model used:} \texttt{qwen3-omni-30b-a3b-instruct} via
  GWDG/SAIA OpenAI-compatible API
  \item \textbf{Input modalities:} normalized requirement text, Vision2Web
  workflow JSON, local resource-path list, route contract, and prototype
  screenshots
  \item \textbf{Message structure:} system message plus one user message
  containing text and prototype images
  \item \textbf{Decoding parameters:} temperature = 0.2, max\_tokens = 20000
  \item \textbf{First use:} 2026-06-04, Europe/Berlin
  \item \textbf{Source:} \texttt{scripts/generate\_track\_b\_ui.py}
  \item \textbf{Result file:}
  \texttt{data/track\_b/items/F10\_gourmania/generated/gwdg\_qwen3\_omni\_30b\_v14\_f10\_smoke/}
\end{itemize}

\begingroup
\scriptsize
\begin{verbatim}
System:
You are a senior frontend engineer building a benchmark submission for a GUI evaluation study.
Create a polished, working website implementation that can be opened directly from a local index.html file.

User:
Prompt ID: TB-GEN-v14

Build a visually grounded, complete single-file HTML implementation for the Track B GUI item below.

Hard requirements:
- Return exactly one complete HTML document ending with </html>.
- Include inline CSS and JavaScript only. No build step, backend, network access, external CDN, or markdown.
- Local assets must be referenced exactly as ../../resources/<subpath> using paths from the resource list.
- Implement every route in the route contract as a <section id="..." data-track-route>.
- Include window.__TRACK_B_ROUTES exactly as given and a showPage(routeId) handler before </body>.
- The first route in window.__TRACK_B_ROUTES must be the default page shown on initial load.

Visual-grounding requirements:
- Use the prototype screenshots for layout, hierarchy, colors, spacing, and visible content priorities.
- Use local image assets whenever they are available. Do not replace provided images with plain text boxes, CSS placeholders, icon-only blocks, or empty cards.
- The homepage must include visible <img> elements when image assets exist.
- Product, recipe, article, service, and video card/list sections should use relevant <img> elements when matching or plausible local images exist.
- Select images by filename and requirement context. For example, filenames that resemble a recipe, cookbook, logo, video thumbnail, or banner should be used for that visible item.
- If no exact image match is obvious, choose the closest local image rather than omitting imagery.
- Reduce prose before removing images, workflow controls, route sections, or JavaScript.

Required workflow-control contract:
- The controls below are machine-checkable requirements. For each item, create an <a> or <button> whose visible text exactly matches visible_text and whose route exactly matches route_id.
- The listed visible_text values are the actual click targets. Context phrases in workflow actions, such as section names, may be headings but do not need to be clickable unless listed below.
- Do not replace visible_text with a synonym, longer phrase, icon-only label, or nearby non-clickable heading.
- Do not implement workflow-required controls as <div>, <span>, <li>, image-only controls, or any other inert element, even if onclick is present.
- If a workflow target is visually a card or image tile, wrap the card content in an <a> or <button> and keep the required visible_text inside that element.
- The required_element example is the safest minimal implementation; copy its text, data-route-target, and onclick route exactly unless the surrounding layout needs <button> instead of <a>.
{workflow_controls}

Implementation guidance:
- Build a recognizable desktop page, not a bare wireframe.
- Prefer a shared header/navigation and shared footer, with route-specific visual content.
- Include enough cards/tables/lists to satisfy the written requirement and workflow validations, but do not reproduce exhaustive long prose.
- Use semantic HTML, visible focus states, alt text for images, and accessible <a>/<button> controls.
- Secondary navigation items may be visible non-functional placeholders only if they are not workflow-required controls.

Route contract:
{route_contract}

Required JavaScript:
<script>
window.__TRACK_B_ROUTES = {route_ids_json};
function showPage(routeId) {
  if (!window.__TRACK_B_ROUTES.includes(routeId)) return;
  document.querySelectorAll('[data-track-route]').forEach(function(section) {
    section.hidden = section.id !== routeId;
  });
  document.querySelectorAll('[data-route-target]').forEach(function(link) {
    link.setAttribute('aria-current', link.dataset.routeTarget === routeId ? 'page' : 'false');
  });
  if (history.replaceState) history.replaceState(null, '', '#' + routeId);
  setTimeout(function() { window.scrollTo(0, 0); }, 0);
}
document.addEventListener('DOMContentLoaded', function() {
  var initial = location.hash ? location.hash.slice(1) : window.__TRACK_B_ROUTES[0];
  showPage(window.__TRACK_B_ROUTES.includes(initial) ? initial : window.__TRACK_B_ROUTES[0]);
});
</script>

Item ID: {item_id}
Source task: {source_task_name}
Source level: {source_level}
Use for dynamic validation: {use_for_dynamic}

Normalized requirement:
{requirement}

Workflow checks:
{workflow}

Available local resource paths:
{resource_paths}

Prototype screenshots are attached after this text in this order:
{prototype_list}
\end{verbatim}
\endgroup

\subsection{TB-GEN-v15: Stricter Route and Click-Target Diagnostic Prompt}
\label{app:prompt-tb-gen-v15}

\textbf{Purpose:} F10 diagnostic prompt introduced after TB-GEN-v14 improved
image usage but still failed static gate checks by using non-contract
\texttt{showPage()} route IDs and replacing the required \texttt{Gefilte Fish}
click target with a generic ``View Recipe'' link.

\textbf{Metadata:}
\begin{itemize}
  \item \textbf{Track / task:} Track B, GUI generation
  \item \textbf{Prompt ID:} TB-GEN-v15
  \item \textbf{Version change:} revises TB-GEN-v14 by forbidding
  non-contract \texttt{showPage()} route IDs and by requiring the clickable
  element's own visible text to exactly match each \texttt{visible\_text}
  \item \textbf{Model used:} \texttt{qwen3-omni-30b-a3b-instruct} via
  GWDG/SAIA OpenAI-compatible API
  \item \textbf{Input modalities:} normalized requirement text, Vision2Web
  workflow JSON, local resource-path list, route contract, and prototype
  screenshots
  \item \textbf{Message structure:} system message plus one user message
  containing text and prototype images
  \item \textbf{Decoding parameters:} temperature = 0.2, max\_tokens = 20000
  \item \textbf{First use:} 2026-06-04, Europe/Berlin
  \item \textbf{Source:} \texttt{scripts/generate\_track\_b\_ui.py}
  \item \textbf{Result file:}
  \texttt{data/track\_b/items/F10\_gourmania/generated/gwdg\_qwen3\_omni\_30b\_v15\_f10\_smoke/}
\end{itemize}

\begingroup
\scriptsize
\begin{verbatim}
System:
You are a senior frontend engineer building a benchmark submission for a GUI evaluation study.
Create a polished, working website implementation that can be opened directly from a local index.html file.

User:
Prompt ID: TB-GEN-v15

Build a visually grounded, complete single-file HTML implementation for the Track B GUI item below.

Hard requirements:
- Return exactly one complete HTML document ending with </html>.
- Include inline CSS and JavaScript only. No build step, backend, network access, external CDN, or markdown.
- Local assets must be referenced exactly as ../../resources/<subpath> using paths from the resource list.
- Implement every route in the route contract as a <section id="..." data-track-route>.
- Include window.__TRACK_B_ROUTES exactly as given and a showPage(routeId) handler before </body>.
- Never call showPage() with a route id that is not listed in window.__TRACK_B_ROUTES.
- Secondary menu items that do not have a route in the route contract must be plain placeholders without onclick, without data-route-target, and without showPage().

Visual-grounding requirements:
- Use the prototype screenshots for layout, hierarchy, colors, spacing, and visible content priorities.
- Use local image assets whenever they are available. Do not replace provided images with plain text boxes, CSS placeholders, icon-only blocks, or empty cards.
- The homepage must include visible <img> elements when image assets exist.
- Product, recipe, article, service, and video card/list sections should use relevant <img> elements when matching or plausible local images exist.
- Select images by filename and requirement context. If no exact image match is obvious, choose the closest local image rather than omitting imagery.
- Reduce prose before removing images, workflow controls, route sections, or JavaScript.

Required workflow-control contract:
- The controls below are machine-checkable requirements. For each item, create an <a> or <button> whose visible text exactly matches visible_text and whose route exactly matches route_id.
- Copy each required_element line into the HTML output exactly once or adapt only the tag name from <a> to <button>. Keep the visible text, data-route-target, and onclick route unchanged.
- The required clickable element's own visible text must be exactly visible_text. Do not use "View", "View Recipe", "Read more", "Open", icons, or longer phrases as the clickable text for required controls.
- The listed visible_text values are the actual click targets. Context phrases in workflow actions, such as section names, may be headings but do not need to be clickable unless listed below.
- Do not implement workflow-required controls as <div>, <span>, <li>, image-only controls, or any other inert element, even if onclick is present.
- If a workflow target is visually a card or image tile, put the required <a> or <button> inside that card, and make the card's clickable label exactly visible_text.
- Do not hide required controls in comments, scripts, templates, off-screen text, display:none, visibility:hidden, or aria-hidden content.
{workflow_controls}

Implementation guidance:
- Build a recognizable desktop page, not a bare wireframe.
- Prefer a shared header/navigation and shared footer, with route-specific visual content.
- Include enough cards/tables/lists to satisfy the written requirement and workflow validations, but do not reproduce exhaustive long prose.
- Use semantic HTML, visible focus states, alt text for images, and accessible <a>/<button> controls.
- Before returning the HTML, self-check these conditions:
  1. Every required_element visible_text appears as visible text inside an <a> or <button>.
  2. Every showPage('...') route is present in window.__TRACK_B_ROUTES and has a matching section id.
  3. The page uses local <img> assets when image assets were provided.

Route contract:
{route_contract}

Required JavaScript:
<script>
window.__TRACK_B_ROUTES = {route_ids_json};
function showPage(routeId) {
  if (!window.__TRACK_B_ROUTES.includes(routeId)) return;
  document.querySelectorAll('[data-track-route]').forEach(function(section) {
    section.hidden = section.id !== routeId;
  });
  document.querySelectorAll('[data-route-target]').forEach(function(link) {
    link.setAttribute('aria-current', link.dataset.routeTarget === routeId ? 'page' : 'false');
  });
  if (history.replaceState) history.replaceState(null, '', '#' + routeId);
  setTimeout(function() { window.scrollTo(0, 0); }, 0);
}
document.addEventListener('DOMContentLoaded', function() {
  var initial = location.hash ? location.hash.slice(1) : window.__TRACK_B_ROUTES[0];
  showPage(window.__TRACK_B_ROUTES.includes(initial) ? initial : window.__TRACK_B_ROUTES[0]);
});
</script>

Item ID: {item_id}
Source task: {source_task_name}
Source level: {source_level}
Use for dynamic validation: {use_for_dynamic}

Normalized requirement:
{requirement}

Workflow checks:
{workflow}

Available local resource paths:
{resource_paths}

Prototype screenshots are attached after this text in this order:
{prototype_list}
\end{verbatim}
\endgroup

\subsection{TB-GEN-v16: Compact-Input Visual Workflow Candidate}
\label{app:prompt-tb-gen-v16}

\textbf{Purpose:} Candidate prompt introduced after the TB-GEN-v15 small-batch
experiment showed that v15 could pass F10 but did not generalize to larger
items because of context/output pressure, repeated workflow input, and exact
text treatment of semantic visual actions such as logo clicks.

\textbf{Metadata:}
\begin{itemize}
  \item \textbf{Track / task:} Track B, GUI generation
  \item \textbf{Prompt ID:} TB-GEN-v16
  \item \textbf{Version change:} revises TB-GEN-v15 by using a compact
  route/evidence workflow summary instead of the full workflow JSON, separating
  exact visible-text controls from semantic/visual targets, adding bounded
  generation limits, and pairing the prompt with deterministic input packing
  (\texttt{compact\_route\_evidence\_v1}, ranked resource paths, and prototype
  image downscaling for compact runs)
  \item \textbf{Model intended for first smoke run:} \texttt{qwen3.6-35b-a3b}
  via GWDG/SAIA OpenAI-compatible API
  \item \textbf{Input modalities:} normalized requirement text, compact
  workflow route/evidence summary, ranked local resource-path list, route
  contract, and prototype screenshots
  \item \textbf{Message structure:} system message plus one user message
  containing text and prototype images
  \item \textbf{Decoding parameters:} planned smoke-test default temperature =
  0.2; max\_tokens chosen per model capability table rather than as a quality
  score
  \item \textbf{First implementation:} 2026-06-05, Europe/Berlin
  \item \textbf{Source:} \texttt{scripts/generate\_track\_b\_ui.py}
  \item \textbf{Result file:} not yet frozen; future smoke runs must record the
  generated directory and \texttt{generation\_metadata.json}
\end{itemize}

\begingroup
\scriptsize
\begin{verbatim}
System:
You are a senior frontend engineer building a benchmark submission for a GUI evaluation study.
Create a polished, working website implementation that can be opened directly from a local index.html file.

User:
Prompt ID: TB-GEN-v16

Build a visually grounded, complete single-file HTML implementation for the Track B GUI item below.

Hard requirements:
- Return exactly one complete HTML document ending with </html>.
- Include inline CSS and JavaScript only. No build step, backend, network access, external CDN, or markdown.
- Local assets must be referenced exactly as ../../resources/<subpath> using paths from the resource list.
- Implement every route in the route contract as a <section id="..." data-track-route>.
- Include window.__TRACK_B_ROUTES exactly as given and a showPage(routeId) handler before </body>.
- Never call showPage() with a route id that is not listed in window.__TRACK_B_ROUTES.
- Secondary menu items that do not have a route in the route contract must be plain placeholders without onclick, without data-route-target, and without showPage().

Visual-grounding requirements:
- Use the prototype screenshots for layout, hierarchy, colors, spacing, and visible content priorities.
- Use local image assets whenever they are available. Do not replace provided images with plain text boxes, CSS placeholders, icon-only blocks, or empty cards.
- Use a bounded number of images: include at least one relevant local image on visual routes when assets exist, but avoid exhaustive galleries.
- Product, recipe, article, service, and video card/list sections should use relevant <img> elements when matching or plausible local images exist.
- Select images by filename and requirement context. If no exact image match is obvious, choose the closest local image from the listed assets rather than omitting imagery.
- Reduce repeated prose before removing images, workflow controls, route sections, validation evidence, or JavaScript.

Required exact workflow-control contract:
- The controls below are machine-checkable exact-text requirements. For each item, create an <a> or <button> whose visible text exactly matches visible_text and whose route exactly matches route_id.
- Copy each required_element line into the HTML output exactly once or adapt only the tag name from <a> to <button>. Keep the visible text, data-route-target, and onclick route unchanged.
- The required clickable element's own visible text must be exactly visible_text. Do not use "View", "View Recipe", "Read more", "Open", icons, or longer phrases as the clickable text for required exact-text controls.
- Do not implement workflow-required exact-text controls as <div>, <span>, <li>, image-only controls, or any other inert element, even if onclick is present.
- Semantic/visual targets listed in the compact workflow summary, such as a logo click, must be clickable and route correctly, but do not render the descriptive phrase literally unless it is actual visible UI text.
{workflow_controls}

Bounded implementation guidance:
- Build a recognizable desktop page, not a bare wireframe.
- Prefer one shared header/navigation and one shared footer, with route-specific visual content.
- For each route, include one clear heading, enough validation evidence text, and at most one compact table/list plus a small card grid when needed.
- Use at most 4 cards per route, 4 rows per table, and 4 bullets per list unless a validation explicitly requires more.
- Use semantic HTML, visible focus states, alt text for images, and accessible <a>/<button> controls.
- Before returning the HTML, self-check these conditions:
  1. Every exact required_element visible_text appears as visible text inside an <a> or <button>.
  2. Every semantic/visual workflow target routes correctly without requiring its descriptive phrase as visible text.
  3. Every showPage('...') route is present in window.__TRACK_B_ROUTES and has a matching section id.
  4. Every route includes the evidence needed by the compact workflow validations.
  5. The page uses local <img> assets when image assets were provided.

Route contract:
{route_contract}

Required JavaScript:
<script>
window.__TRACK_B_ROUTES = {route_ids_json};
function showPage(routeId) {
  if (!window.__TRACK_B_ROUTES.includes(routeId)) return;
  document.querySelectorAll('[data-track-route]').forEach(function(section) {
    section.hidden = section.id !== routeId;
  });
  document.querySelectorAll('[data-route-target]').forEach(function(link) {
    link.setAttribute('aria-current', link.dataset.routeTarget === routeId ? 'page' : 'false');
  });
  if (history.replaceState) history.replaceState(null, '', '#' + routeId);
  setTimeout(function() { window.scrollTo(0, 0); }, 0);
}
document.addEventListener('DOMContentLoaded', function() {
  var initial = location.hash ? location.hash.slice(1) : window.__TRACK_B_ROUTES[0];
  showPage(window.__TRACK_B_ROUTES.includes(initial) ? initial : window.__TRACK_B_ROUTES[0]);
});
</script>

Item ID: {item_id}
Source task: {source_task_name}
Source level: {source_level}
Use for dynamic validation: {use_for_dynamic}

Normalized requirement:
{requirement}

Compact workflow checks, deduplicated from workflow.json:
{workflow}

Relevant local resource paths:
{resource_paths}

Prototype screenshots are attached after this text in this order:
{prototype_list}
\end{verbatim}
\endgroup

\section{Track B Static-Visual Judge and Human Review Prompts}
\label{app:track-b-visual-judge-prompts}

\subsection{LB-JUDGE-v1: Base Static-Visual Checklist Prompt}
\label{app:prompt-lb-judge-v1-base}

\textbf{Purpose:} Static visual scoring for standardized Track~B page
screenshots. The prompt asks a multimodal LLM to answer the same 16 binary
checklist items used by the human review interface, returning machine-readable
JSON.

\textbf{Metadata:}
\begin{itemize}
  \item \textbf{Track / task:} Track B, static visual judging
  \item \textbf{Prompt ID:} LB-JUDGE-v1-base
  \item \textbf{Models used in development comparisons:} \texttt{internvl3.5-30b-a3b},
  \texttt{qwen3-omni-30b-a3b-instruct}, \texttt{gemma-4-31b-it},
  \texttt{gpt-4o-mini}, and \texttt{gpt-4.1-mini}, depending on provider
  availability
  \item \textbf{Input modalities:} one standardized PNG screenshot plus text
  checklist
  \item \textbf{Message structure:} system message plus one user message with
  text and image
  \item \textbf{Decoding parameters:} temperature = 0.0, max\_tokens = 800
  \item \textbf{Source:} \texttt{scripts/run\_visual\_judge.py}
  \item \textbf{Result files:}
  \texttt{data/track\_b/visual\_human\_review/llm\_judge\_*.json} and
  \texttt{human\_vs\_llm\_comparison*.json}
\end{itemize}

\begingroup
\scriptsize
\begin{verbatim}
System:
You are a careful visual UI reviewer. You see ONE screenshot of ONE page of
a web interface. Judge ONLY what is visible. Do not guess about interactivity
or pages you cannot see. Answer each checklist item strictly true or false.
Return ONLY valid JSON, no prose.

User:
Answer every checklist item for this screenshot (true = page satisfies it).

Layout & Visual Hierarchy
 L1: There is a clear primary element or focal point.
 L2: Elements are free of overlap or obvious misalignment.
 L3: Alignment follows a consistent grid/structure.
 L4: Spacing and padding are consistent across the page.
Information Organization & Clarity
 O1: Related content is grouped into clear sections.
 O2: The page is uncluttered, not visually overwhelming.
 O3: Sections are labelled or easy to tell apart.
 O4: The main purpose/content is clear at a glance.
Typography & Readability
 T1: Clear text hierarchy (headings vs body distinguishable).
 T2: Body text is a comfortably readable size.
 T3: Line length / text width is reasonable.
 T4: Font usage is consistent (no clashing typefaces).
Visual Consistency
 C1: Buttons and controls share a consistent style.
 C2: Colour palette is coherent (no clashing colours).
 C3: Navigation / header styling is consistent.
 C4: Repeated components look uniform.

Return JSON exactly:
{"answers": {"L1": true, "L2": true, "L3": true, "L4": true,
 "O1": true, "O2": true, "O3": true, "O4": true,
 "T1": true, "T2": true, "T3": true, "T4": true,
 "C1": true, "C2": true, "C3": true, "C4": true}, "confidence": 0.0}
\end{verbatim}
\endgroup

\subsection{LB-JUDGE-v1: Strict Variant}
\label{app:prompt-lb-judge-v1-strict}

\textbf{Purpose:} Bias-reduction variant introduced after base prompts across
several vision models produced all-true or near-all-true outputs. The checklist
items are identical to LB-JUDGE-v1-base; only the system instruction and the
lead-in before the JSON schema change.

\textbf{Metadata:}
\begin{itemize}
  \item \textbf{Track / task:} Track B, static visual judging
  \item \textbf{Prompt ID:} LB-JUDGE-v1-strict
  \item \textbf{Version change:} adds fault-finding instruction and changes
  ``Return JSON exactly'' to ``Look for problems first. Return JSON exactly''
  \item \textbf{Source:} \texttt{scripts/run\_visual\_judge.py}
\end{itemize}

\begingroup
\scriptsize
\begin{verbatim}
System:
You are a STRICT, critical visual UI reviewer evaluating an AI-generated web
page. Most AI-generated pages contain at least some real visual flaws
(inconsistent spacing, weak contrast, misalignment, cramped or oversized
text, clashing colours, uneven components). Your job is to FIND those flaws,
not to approve. Mark an item true ONLY if the criterion is clearly and fully
satisfied with no visible issue. If there is ANY visible problem, or you are
not clearly confident, mark it false. A page that is true on every item is
rare; before answering, identify the single worst visual problem on the page
and make sure the relevant item is false. Return ONLY valid JSON, no prose.

User:
Same 16-item checklist as LB-JUDGE-v1-base, with this lead-in before the
schema:

Look for problems first. Return JSON exactly:
\end{verbatim}
\endgroup

\subsection{LB-JUDGE-v1: Strict Few-Shot Variant}
\label{app:prompt-lb-judge-v1-strict-fewshot}

\textbf{Purpose:} Calibration variant for static visual scoring. It uses the
strict system and user prompts above, then prepends three human-labelled
example screenshots before the scored screenshot. The example pages are
excluded from evaluation when computing human-vs-LLM agreement.

\textbf{Metadata:}
\begin{itemize}
  \item \textbf{Track / task:} Track B, static visual judging with calibration
  examples
  \item \textbf{Prompt ID:} LB-JUDGE-v1-strict-fewshot
  \item \textbf{Example pages:} \texttt{F10\_gourmania::01\_homepage}
  (human page score 0.50), \texttt{F01\_1daycloud::04\_contact}
  (0.75), and \texttt{F01\_1daycloud::02\_academy} (1.00)
  \item \textbf{Message structure:} strict system message; for each example,
  one user image+checklist message followed by one assistant JSON answer; then
  the scored user image+checklist message
  \item \textbf{Source:} \texttt{scripts/run\_visual\_judge.py}
\end{itemize}

\begingroup
\scriptsize
\begin{verbatim}
System:
Use the LB-JUDGE-v1-strict system message.

Few-shot messages:
For each calibration page, send the same LB-JUDGE-v1-strict user checklist
with the example screenshot, followed by an assistant message:

{"answers": {"L1": <human label>, ..., "C4": <human label>},
 "confidence": 1.0}

Scored message:
Send the same LB-JUDGE-v1-strict user checklist with the target screenshot.
\end{verbatim}
\endgroup

\subsection{HUMAN-VISUAL-v1: Human Review Checklist}
\label{app:prompt-human-visual-v1}

\textbf{Purpose:} Human annotation protocol for validating the static visual
checklist against LLM judges. The self-contained review page shows one
standardized screenshot at a time and asks the rater to answer the same
16 binary checklist items as LB-JUDGE-v1.

\textbf{Metadata:}
\begin{itemize}
  \item \textbf{Track / task:} Track B, human static-visual annotation
  \item \textbf{Prompt ID:} HUMAN-VISUAL-v1
  \item \textbf{Interface source:} \texttt{scripts/build\_visual\_human\_review.py}
  \item \textbf{Original export:}
  \texttt{data/track\_b/visual\_human\_review/visual\_human\_review.json}
  \item \textbf{Extended export:}
  \texttt{data/track\_b/visual\_human\_review/visual\_human\_review\_extended.json}
  \item \textbf{Extended setting:} deterministic shuffled order, run names
  hidden from the page, separate localStorage key and export name
\end{itemize}

\begingroup
\scriptsize
\begin{verbatim}
Static-Visual Human Review

For each screenshot, answer the same 16 checklist items:
L1-L4 Layout & Visual Hierarchy
O1-O4 Information Organization & Clarity
T1-T4 Typography & Readability
C1-C4 Visual Consistency

Each item is answered Yes or No. The interface also provides All Yes and
All No buttons for a page; the rater can then correct individual items before
exporting JSON.
\end{verbatim}
\endgroup

\section{Prompt Slots for Later Experiments}
\label{app:future-prompt-slots}

Future thesis-impacting prompts should be added here before or immediately after
their first experimental use. Expected prompt families include Track B static
rubric scoring prompts, LLM-based dynamic validation prompts if such validators
are later used, output-schema prompts, and human annotation instructions.
